% =========================================================
\chapter{Métaux et alliages 1 : l'aluminium, le cuivre et leurs alliages}
% =========================================================

\begin{rembox}[title=En bref]
\textbf{De quoi parle ce chapitre ?} Les métaux, à travers deux exemples : l'\textbf{aluminium}
et le \textbf{cuivre}. On voit d'où ils viennent (ressources, énergie, recyclage, extraction
guidée par le diagramme d'Ellingham), puis le c\oe ur du chapitre : comment la
\textbf{microstructure} se forme (solidification, ségrégation, maturation, précipitation) et
comment on l'utilise pour \textbf{durcir} un métal. \\
\textbf{Les quatre mécanismes de durcissement vus ici} : écrouissage (défait par le recuit),
solution solide, durcissement structural (précipités), renfort par une seconde phase dure. Tous
(sauf le dernier) consistent à mettre des \textbf{obstacles sur le chemin des dislocations}.
\end{rembox}

\section{Panorama des métaux : propriétés, ressources, énergie}

\begin{propbox}[title=Ce qui distingue les métaux et alliages]
Un élément principal (Fe, Al, Cu, Mg, Ti, Ni\ldots) plus des \textbf{éléments d'alliage}. Moins
rigides et moins durs, souvent plus lourds que les céramiques, mais \textbf{ductiles} et
\textbf{tenaces} : c'est ce qui en fait les matériaux de structure par excellence. Ils peuvent
s'oxyder, ont une empreinte énergie/carbone non négligeable, mais se recyclent facilement.
\end{propbox}

\begin{rembox}[title={Une coïncidence utile}]
Fe, Al, Ti, Mg et Ni ont presque le \textbf{même $E/\rho$} ($\approx 25$\,GPa/(Mg/m$^3$)) : en
traction, à rigidité égale, une pièce en aluminium pèse autant qu'une pièce en acier. L'avantage
des alliages légers n'apparaît qu'en flexion ($E/\rho^2$, chapitre~2) ou en résistance.
\end{rembox}

\begin{defnbox}[title={Réserves et ressources}]
\textbf{Réserves} : gisements connus exploitables \textbf{avec profit}, aujourd'hui, avec les
techniques d'aujourd'hui (elles augmentent quand le prix monte ou que la technique progresse).
\textbf{Ressources} : tout ce dont on peut envisager l'extraction, l'autre extrême.
\end{defnbox}

Le coût énergétique d'un métal augmente avec sa \textbf{dilution} dans l'écorce terrestre et
avec la stabilité de son oxyde : l'aluminium, très abondant mais très lié à l'oxygène, coûte
cher en énergie ; l'or, peu lié mais très dilué, encore plus.

\subsection*{Le recyclage est limité par la croissance de la demande}

\slide[0.6]{c03a-38}{Ce qu'on recycle aujourd'hui a été produit il y a $n$ années, quand on produisait moins.}

\begin{thmbox}[title=Taux de recyclage maximal]
Si la production croît de $i$ par an et qu'un objet dure $n$ ans en moyenne, le métal récupérable
aujourd'hui vient de la production d'il y a $n$ ans ($A$), plus faible que celle d'aujourd'hui
($B$) :
\[ F \le \frac{A}{B} = (1+i)^{-n} \]
\end{thmbox}

\begin{exbox}[title={Exemple chiffré (pas dans les diapositives)}]
Aluminium : $n = 15$ ans, croissance $i = 4\,\%$ par an :
$F \le 1{,}04^{-15} = 0{,}56$. Même en recyclant \emph{tout}, on ne couvrirait que 56\,\% de la
production : le reste doit venir du minerai.
\end{exbox}

\begin{goatbox}
Les métaux sont les matériaux qui plient au lieu de casser. On en a plein sous nos pieds, mais
les extraire coûte de l'énergie, surtout quand le métal est très attaché à l'oxygène (aluminium)
ou très dispersé (or). Recycler aide beaucoup, mais tant qu'on en veut chaque année plus, les
vieux objets à recycler ne suffisent pas à tout fournir.
\end{goatbox}

\section{Produire l'aluminium et le cuivre}

\subsection*{Le diagramme d'Ellingham}

\slide[0.7]{c03a-45}{Diagramme d'Ellingham : plus la droite d'un oxyde est basse, plus l'oxyde est stable. Al$_2$O$_3$ est tout en bas.}

\begin{defnbox}[title={Lire le diagramme d'Ellingham}]
Pour chaque réaction \og métal + O$_2$ $\to$ oxyde\fg{} (rapportée à une mole d'O$_2$), on trace
$\Delta G^\circ$ en fonction de $T$.
\begin{itemize}
  \item \textbf{Plus la droite est basse} ($\Delta G^\circ$ très négatif), \textbf{plus
        l'oxyde est stable}, plus il est difficile d'en extraire le métal.
  \item Un élément dont la droite est \textbf{en dessous} de celle d'un autre peut
        \textbf{réduire} l'oxyde de cet autre (il lui \og vole\fg{} l'oxygène).
  \item La droite C + O$_2$ $\to$ 2CO \textbf{descend} quand $T$ augmente : à haute
        température, le carbone devient capable de réduire de nombreux oxydes (fer, cuivre),
        mais pas Al$_2$O$_3$ dans des conditions raisonnables.
\end{itemize}
(Une simplification : les droites dépendent aussi des pressions partielles des gaz.)
\end{defnbox}

\subsection*{Le cuivre : facile à réduire, difficile à trouver}

Le cuivre ne représente que 50\,ppm de l'écorce : le minerai (sulfures, souvent
CuFeS$_2$) est pauvre et coûteux à concentrer. En revanche, Cu est presque \textbf{noble} :
réduire ses composés demande peu d'énergie. La voie \textbf{pyrométallurgique} : concentration
par flottation, fusion en \textbf{matte} (Cu$_2$S--FeS), \textbf{convertisseur} où l'on souffle de
l'air (le fer part dans la scorie avec la silice, le soufre en SO$_2$), cuivre blister, puis
\textbf{raffinage électrolytique} (anodes impures, cathodes de cuivre pur à 99,99\,\%). Il existe
aussi une voie \textbf{hydrométallurgique} (lixiviation à l'acide sulfurique, puis
électrolyse).

\subsection*{L'aluminium : abondant mais très lié à l'oxygène}

Son oxyde est si stable qu'on ne peut pas le réduire au carbone : il faut
l'\textbf{électrolyse}.
\begin{enumerate}
  \item \textbf{Procédé Bayer} : de la bauxite on tire de l'alumine Al$_2$O$_3$ pure.
  \item \textbf{Procédé Hall--Héroult} (1886) : l'alumine est dissoute dans de la cryolithe
        fondue ($\approx 960$\,°C) et électrolysée : $2\,\text{Al}^{3+} + 6e^- \to 2\,\text{Al}$
        à la cathode ; à l'anode en carbone, $3\,\text{O}^{2-} \to \frac32\text{O}_2 + 6e^-$,
        l'oxygène brûlant l'anode en CO$_2$.
\end{enumerate}

\slide[0.6]{c03a-58}{La cuve d'électrolyse à anodes précuites.}

\begin{propbox}[title={L'énergie de l'aluminium}]
Production mondiale $\approx 100$\,Mt/an, dont $\approx 30\,\%$ recyclé. Produire 1\,kg depuis le
minerai coûte $12 \pm 2$\,kWh (limite théorique 6,5\,kWh/kg) ; le recycler, $\approx 1$\,kWh/kg
(10\,\%).
\end{propbox}

\begin{exbox}[title=Le petit exercice des diapositives]
À combien d'accélérations de 0 à 100\,km/h correspondent 12\,kWh/kg ? Par kilogramme,
$\frac12 v^2 = \frac12(10^5/3600)^2 = 386$\,J ; $12$\,kWh $= 43\,200$\,kJ ; donc
$43\,200\,000/386 \approx 112\,000$ accélérations, soit 31 par jour pendant 10 ans. Alléger un
véhicule avec de l'aluminium peut donc rembourser l'énergie de sa production --- à évaluer au
cas par cas.
\end{exbox}

\begin{propbox}[title={Pourquoi l'aluminium ne rouille pas}]
Son oxyde, très stable, forme une couche dure, inerte et \textbf{mécaniquement compatible} avec
le métal : le rapport de \textbf{Pilling--Bedworth} (volume de l'oxyde / volume du métal
consommé) est proche de 1, donc la couche se forme sans fissures ni décollement et protège le
métal. On dit que l'aluminium est \textbf{passivé}.
\end{propbox}

\begin{rembox}[title=Mise en forme]
Fusion, puis soit coulée en lingot suivie de \textbf{laminage} (tôles, feuilles),
\textbf{extrusion} (profilés, tubes) ou \textbf{forgeage} ; soit \textbf{moulage} direct
(sable, cire perdue, moule permanent, sous pression) ; et désormais l'impression 3D (fusion
laser de poudre).
\end{rembox}

\begin{goatbox}
Le diagramme d'Ellingham est un classement de \og qui tient le plus fort à l'oxygène\fg{}.
L'aluminium y tient tellement qu'aucun charbon ne peut le lui arracher : il faut l'électricité,
et beaucoup. Le cuivre, lui, lâche son oxygène (ou son soufre) facilement, mais il est rare et
dispersé. La même passion de l'aluminium pour l'oxygène le protège : une fine peau d'oxyde,
bien collée, l'empêche de rouiller.
\end{goatbox}
